\documentclass[11pt]{article}
\usepackage[a4paper,margin=27mm]{geometry}
\usepackage{amsmath,amssymb,amsthm,hyperref}
\newcommand{\ddc}{dd^c}
\newcommand{\NA}{\mathrm{NA}}
\newcommand{\MA}{\operatorname{MA}}
\newcommand{\Ent}{\operatorname{Ent}}
\newtheorem{theorem}{Theorem}
\title{An actual projective-line ray calibrating energy and entropy slopes}
\author{}
\date{6 October 2026}
\begin{document}
\maketitle

This is an unreviewed, explicit calculation for a known projective model.
It supplies a normalization check for the residual P19.1 research, rather
than a solution of the arbitrary compact K\"ahler maximal-ray problem.
Frozen native inputs131 and132 are unchanged.

\section{The ray, the current convention, and time}

Put $\ddc=\frac{i}{2\pi}\partial\bar\partial$ and use
$t=-\log|\zeta|^2$. Thus $\ddc\log|f|^2$ is the current of the
divisor of a holomorphic function $f$. On $\mathbf P^1$ take
\[
 \omega=\ddc\log(1+|z|^2),\qquad \int_{\mathbf P^1}\omega=1.
\]
For $t\geq0$ define
\[
 u_t(z)=\log\frac{|z|^2+e^{-t}}{1+|z|^2}.
 \tag{1}
\]
In the coordinate $w=1/z$ the same function is
$\log((1+e^{-t}|w|^2)/(1+|w|^2))$, so it is smooth on the
whole projective line for every finite time. It satisfies
$-t\leq u_t\leq0$ and $u_0=0$.
Use the probability measure $\MA_\omega(u)=\omega+\ddc u$ and
\[
 E_\omega(u)=\frac12\int_{\mathbf P^1}u\,
                  (\omega+\omega_u),\qquad
 H_\omega(u)=\Ent(\MA_\omega(u)\mid\omega).
 \tag{2}
\]

\begin{theorem}
The path (1) is an actual bounded weak geodesic ray from zero. Its
normalized energy is $E_\omega(u_t)=-t/2$. Its usual non-Archimedean
datum is the model function
$\phi(v)=-\min\{v(\mathcal I_0),1\}$, where $\mathcal I_0$ is the
ideal of $0\in\mathbf P^1$. Its non-Archimedean probability measure
is $\delta_{\operatorname{ord}_0}$, and $H^{\NA}(\phi)=1$.
Moreover,
\[
 H_\omega(u_t)
 =t-2+\frac{2t e^{-t}}{1-e^{-t}}\quad(t>0),\qquad
 \lim_{t\to\infty}\frac{H_\omega(u_t)}t=1.
 \tag{3}
\]
\end{theorem}

\begin{proof}
The local total potential, after adding the local background potential,
is
\[
 \log(1+|z|^2)+u_{-\log|\zeta|^2}(z)
       =\log(|z|^2+|\zeta|^2).
 \tag{4}
\]
It is psh. On a time annulus $\zeta\ne0$ its complex Hessian has
rank one: the Hessian of $\log\|Z\|^2$ on $\mathbf C^2\setminus
\{0\}$ has the radial complex line as kernel and one positive
eigenvalue. Its square therefore vanishes. In the other projective
coordinate the total background potential is
$\log(1+|\zeta|^2|w|^2)$; its Hessian also has rank at most one.
Thus (4) defines a positive global total current with vanishing
homogeneous Monge--Amp\`ere measure on each annulus.

The total potentials are smooth and bounded on every closed finite
annulus times $\mathbf P^1$. The usual bounded Bedford--Taylor
comparison principle identifies this solution with the weak
geodesic Perron envelope with its boundary values. The same argument
on every pair of finite times proves the asserted weak ray.
Only this standard bounded Dirichlet characterization is imported
for the geodesic terminology.

Write $a=e^{-t}$ and $y=|z|^2$. Integrating out the angular
variable, the reference and current probability measures are
\[
 d\omega=\frac{dy}{(1+y)^2},\qquad
 d\MA_\omega(u_t)=\frac{a\,dy}{(y+a)^2},\qquad y\in[0,\infty).
 \tag{5}
\]
Both integrals equal one. Since
$\dot u_t=-a/(y+a)$, the smooth energy derivative gives
\[
 \frac{d}{dt}E_\omega(u_t)
 =\int\dot u_t\,d\MA_\omega(u_t)
 =-\int_0^\infty\frac{a^2\,dy}{(y+a)^3}
 =-\frac12.
 \tag{6}
\]
Here the derivative identity also follows directly by differentiating
(2) and integrating by parts on the compact curve.
The energy is zero at time zero, proving $E_\omega(u_t)=-t/2$.

The Radon--Nikodym density in (5) is
$\rho_t=a(1+y)^2/(y+a)^2$. Substitute $y=ax$ in its entropy:
\[
 H_\omega(u_t)
 =t+2I(a)-2I(1),\qquad
 I(a)=\int_0^\infty\frac{\log(1+ax)}{(1+x)^2}\,dx.
 \tag{7}
\]
The boundary term in integration by parts is zero, so for
$0<a<1$,
\[
 I(a)=a\int_0^\infty\frac{dx}{(1+x)(1+ax)}
      =-\frac{a\log a}{1-a}.
 \tag{8}
\]
Direct integration gives $I(1)=1$. Equations (7)--(8) give (3).
They also give $H_\omega(u_t)\to0$ as $t\downarrow0$; its infinite-time
slope is exactly one.

For completeness compute the actual datum and its model measure.
Blow up $\mathbf P^1\times\mathbf P^1$ at $(z,\zeta)=(0,0)$,
and denote the exceptional divisor by $F$ and the strict transform
of the central fiber by $C$. The current in (4) has generic
coefficient one along $F$ and coefficient zero along $C$.
The central multiplicities are $b_F=b_C=1$.
The coherent ideal $(z,\zeta)$ consequently gives the vertical
divisor $D=-F$ and model function
\[
 f_D(v)=-\min\{v(\mathcal I_0),1\}.
 \tag{9}
\]
The model interpretation uses the usual normalized Gauss extension,
whose time weight is one. Thus (9) is obtained from the total
potential, rather than assigned to the path.

Let $\alpha$ be the pullback of $[\omega]$ from the first factor.
The relatively nef model class is $\alpha-F$.
The intersection numbers are
\[
 \alpha\cdot F=0,\quad F^2=-1,\quad
 \alpha\cdot C=1,\quad F\cdot C=1.
 \tag{10}
\]
Hence $(\alpha-F)\cdot F=1$ and
$(\alpha-F)\cdot C=0$. The defining model measure
$\sum_Z b_Z((\alpha-F)\cdot Z)\delta_{v_Z}$ is therefore
$\delta_{\operatorname{ord}_0}$. Discrepancy over the smooth
curve, with $A_{\mathbf P^1}(\operatorname{ord}_0)=1$, gives
$H^{\NA}(\phi)=1$.
Finally $(\alpha-F)^2=-1$, so its model energy
$(\alpha-F)^2/2=-1/2$ agrees with (6).
These are relatively nef model calculations; no absolute bigness
is asserted. If an absolutely big representative is desired,
one may add a central-fiber constant and subtract that same
affine constant from the ray afterwards.
\end{proof}

\section{The factor under the other time convention}

Keep the same current convention and cohomological model, but
write $s=-\log|\zeta|$. Then the same total potential has ray
$\widetilde u_s=u_{2s}$. Equations (3) and (6) imply
\[
 E_\omega(\widetilde u_s)=-s,\qquad
 \lim_{s\to\infty}\frac{H_\omega(\widetilde u_s)}s=2.
 \tag{11}
\]
The model energy and entropy remain $-1/2$ and $1$.
Consequently the two equalities
$E_\omega(\widetilde u_s)=sE^{\NA}(\phi)$ and
$H_\omega^\infty(\widetilde u)=H^{\NA}(\phi)$ would both be false
with these unchanged conventions.

There is an equally coherent convention with
$\widehat{\ddc}=\frac{i}{\pi}\partial\bar\partial$,
the same background form $\omega$, and time $s=-\log|\zeta|$.
Set $\widehat u_s=u_{2s}/2$. Then
$\omega+\widehat{\ddc}\widehat u_s=\omega+\ddc u_{2s}$,
so its probability measures and entropy are identical, whereas
its normalized energy is halved. Its total divisor coefficients
are again one because
$\widehat{\ddc}\log|f|=[f=0]$. Thus
\[
 E_\omega^{\,\widehat{\ddc}}(\widehat u_s)=-s/2,\qquad
 H_\omega^\infty(\widehat u)=2H^{\NA}(\phi).
 \tag{12}
\]
This explicit nonzero-entropy model distinguishes the conventions;
the zero metric alone cannot detect this factor.

\section{What this calculation settles}

With $\ddc=i\partial\bar\partial/(2\pi)$ and log-squared divisor
coefficients, the time parameter $-\log|\zeta|^2$ makes both
the energy and probability-entropy slopes equal to their
cohomological model values in this example.
The provisional continuation178 instead printed
$-\log|\zeta|$ with those same other conventions.
Equations (11) provide a direct mathematical correction to that
combination. A full corrected transcendental proof must transport
the ray correspondence, mixed slopes, Mabuchi formula and entropy
lower bound consistently; this calibration does not itself
establish those general transfers.

The construction is a smooth polarized projective model, already
within the known state of the art. It does not close original P19.1
or any other original atom, and has received no independent PASS.

\end{document}
